\documentclass{article}

\usepackage{amsmath,amssymb}
\usepackage{xurl}
\usepackage[hidelinks]{hyperref}
\setlength{\emergencystretch}{2em}

% Shared commands for notes in docs/paper-gaps/.

\DeclareUnicodeCharacter{2080}{\ifmmode{}_{0}\else\textsubscript{0}\fi}
\DeclareUnicodeCharacter{2081}{\ifmmode{}_{1}\else\textsubscript{1}\fi}
\DeclareUnicodeCharacter{2082}{\ifmmode{}_{2}\else\textsubscript{2}\fi}
\DeclareUnicodeCharacter{2099}{\ifmmode{}_{n}\else\textsubscript{n}\fi}

\newcommand{\Lean}{\textsc{Lean}}
\ExplSyntaxOn
\NewDocumentCommand{\leanid}{m}{
  \ifmmode
    \textsf{\detokenize{#1}}
  \else
    \group_begin:
    \urlstyle{sf}
    \tl_set:Nn \l_tmpa_tl {#1}
    \tl_replace_all:Nnn \l_tmpa_tl {\_} {_}
    \exp_args:NV \path \l_tmpa_tl
    \group_end:
  \fi
}
\ExplSyntaxOff
\providecommand{\tr}{\operatorname{tr}}
\newcommand{\tnleanIssueBase}{https://github.com/LionSR/TNLean/issues/}
\newcommand{\tnleanPrBase}{https://github.com/LionSR/TNLean/pull/}
\newcommand{\tnleanDocBase}{https://sirui-lu.com/TNLean/docs/}
\newcommand{\ghissue}[1]{\href{\tnleanIssueBase#1}{GitHub issue \##1}}
\newcommand{\ghpr}[1]{\href{\tnleanPrBase#1}{PR \##1}}
\newcommand{\doclink}[1]{\href{\tnleanDocBase#1.html}{\path{#1}}}

% Dirac notation and the identity operator, as in blueprint/src/macros/common.tex.
% \providecommand keeps note-local definitions authoritative.
\usepackage{dsfont}
\providecommand{\ket}[1]{|#1\rangle}
\providecommand{\bra}[1]{\langle #1|}
\providecommand{\braket}[2]{\langle #1 | #2 \rangle}
\providecommand{\ketbra}[2]{|#1\rangle\!\langle #2|}
\providecommand{\Id}{\mathds{1}}
\providecommand{\one}{\mathds{1}}
\providecommand{\C}{\mathbb{C}}
\providecommand{\MN}[1]{M_{#1}(\C)}
\providecommand{\MD}{M_D(\C)}

% External referencing for paper sources.  The build helper writes sanitized
% auxiliary files under build/paper-aux/ and excludes duplicated source labels.
\usepackage{xr}

% Import a generated paper auxiliary file with a caller-supplied label prefix.
% The candidate paths support builds from docs/paper-gaps/, the repository root,
% and build/paper-aux/ itself.
\newcommand{\importpaperaux}[2]{%
  \IfFileExists{../../build/paper-aux/#2.aux}{%
    \externaldocument[#1]{../../build/paper-aux/#2}%
  }{%
    \IfFileExists{build/paper-aux/#2.aux}{%
      \externaldocument[#1]{build/paper-aux/#2}%
    }{%
      \IfFileExists{#2.aux}{%
        \externaldocument[#1]{#2}%
      }{}%
    }%
  }%
}

% General external reference macros
\newcommand{\paperref}[2]{\ref{#1-#2}}
\newcommand{\papereqref}[2]{(\ref{#1-#2})}

% CPSV16 source (1606.00608 / MPDO-22-12-17-2.tex) external document and shortcuts
\importpaperaux{cpsv16-}{1606.00608/MPDO-22-12-17-2-tnlean}
\newcommand{\cpsvref}[1]{\paperref{cpsv16}{#1}}
\newcommand{\cpsveqref}[1]{\papereqref{cpsv16}{#1}}

% Verdict marker: \gapnote{<kind>}{<status>}. Kind is the policy
% classification (clarification, local-correction, scope-restriction,
% unfaithful, false-source, open-gap); status is open, wip (elimination
% underway), resolved, or historical. Machine-read by the site index;
% prints a verdict line.
\newcommand{\gapnote}[2]{\par\noindent\textbf{Verdict:} #1 (#2).\par}


\title{The Direct-Sum Input from the Older MPS Representation Paper}
\author{}
\date{2026-05-05}

\begin{document}
\maketitle

\section{Source statement}

The finite-length direct-sum input used behind the block-injective
canonical-form route is not proved in the 2016 MPDO source. The 2016 paper
cites the older representation paper for the block-injective-after-blocking
step \cite[lines~340--345]{Cirac2016MPDO_arXiv}. The local source for the
older argument is \path{Papers/quant-ph_0608197/MPSarchive.tex}.

That source first introduces the injectivity condition $C1$: for some
$L_0$, the map
\[
  \Gamma_{L_0}(X) = \sum_{i_1,\ldots,i_{L_0}}
      \operatorname{tr}(X A_{i_1}\cdots A_{i_{L_0}})
      |i_1,\ldots,i_{L_0}\rangle
\]
is injective. It immediately identifies this with the exact-length word
products of length $L_0$ spanning the full matrix algebra
\cite[lines~888--908]{PerezGarcia2007MPSReps}.

The technical algebraic lemma then says that, for nonzero $D\times D$
matrices $C$, every matrix $X$ is a finite sum
\[
  X = \sum_i R_i C S_i.
\]
The proof reduces $C$ to a matrix unit by elementary matrix operations and
then reconstructs $X$ from permuted copies of that matrix unit
\cite[lines~1333--1344]{PerezGarcia2007MPSReps}.

The next lemma is the direct-sum statement. Its local hypotheses are important:
the tensor has $b \ge 2$ canonical blocks $A^1,\ldots,A^b$, each block
satisfies $C1$, the common length is
$L_0=\max_j L_0^{A^j}$, the block sizes are ordered, and the block states are
assumed pairwise different
\cite[lines~1321--1330]{PerezGarcia2007MPSReps}. Under these hypotheses, if
$L \ge 3(b-1)(L_0+1)$, then
\[
  \sum_{j=1}^b \mathcal G_L^{A^j}
\]
is a direct sum \cite[lines~1346--1408]{PerezGarcia2007MPSReps}. The proof is
by induction on the number of blocks. The base case uses $C1$ and the
matrix-factorization lemma to show that a nontrivial linear dependence would
force inclusion, then equality, of the two length-$L_0+1$ MPS subspaces,
contradicting uniqueness of the associated injective parent-space
characterization. The induction step isolates one block with a local
projector, applies the induction hypothesis to the remaining blocks, and then
uses $C1$ plus the factorization lemma to force the remaining coefficient to
vanish.

\section{MPS translation}

For the non-periodic Fundamental Theorem route, the mathematical consequence is
homogeneous and positive-length. After the canonical-form and
basis-of-normal-tensors reduction, the selected representatives should admit one
length at which their simultaneous word evaluations separate the direct-sum
blocks. In trace-dual form, if
\[
  \sum_{j=1}^g \operatorname{tr}(\Delta_j A_j^w)=0
  \qquad\text{for every word } w \text{ of length } L,
\]
then each $\Delta_j$ is zero. This is the finite-length direct-sum span input
needed by the BNT pair-separation route.

This statement is not a consequence of the empty word and is not a cumulative
span statement. The exact length $L$ is part of the content. Therefore the
formal proof should not replace the source input by a hypothesis that only says
the cumulative pair algebra is full, nor by a theorem that uses the length-zero
identity as padding.

\section{Formal boundary}

The matrix-factorization content of Lemma~\texttt{lem1} from the source paper
is now formalized in two steps.\footnote{%
\leanid{Matrix.span_range_mul_nonzero_mul_eq_top} for the abstract
matrix-algebra form; \leanid{MPSTensor.span_range_evalWord_mul_nonzero_mul_evalWord_eq_top}
for the word-product translation.}
The first says that, for any nonzero matrix $C$, the set of products
$\{RCS : R,S\in M_D(\mathbb C)\}$ spans $M_D(\mathbb C)$. The second
translates this into the relevant word-product form: if a tensor is $L$-block
injective and $X\ne0$, then
$\operatorname{span}\{A_u X A_v : |u|=|v|=L\}=M_D(\mathbb C)$.

The next part of the two-block direct-sum argument is also formalized. The
three-block trace relation forces, in the size-ordering convention $D_B\le D_A$,
an inclusion $\mathcal G_L^A\subseteq\mathcal G_L^B$ of length-$L$ image
spaces. Since block injectivity identifies $\dim\mathcal G_L^A=D_A^2$ and
$\dim\mathcal G_L^B=D_B^2$, the inclusion and the size ordering together
force $D_A=D_B$ and $\mathcal G_L^A=\mathcal G_L^B$.\footnote{%
\leanid{MPSTensor.bondDim_eq_and_groundSpace_eq_of_three_block_trace_relation_left_of_dim_ge}.}
The strict-size branch - if $D_B<D_A$, no nonzero coefficient on the larger
block can satisfy the homogeneous three-block trace relation - is formalized
separately,\footnote{%
\leanid{MPSTensor.not_three_block_trace_relation_left_of_isNBlkInjective_of_dim_gt};
\leanid{MPSTensor.pairTraceSeparatingAt_threeBlock_of_isNBlkInjective_of_dim_gt}
for the trace-separation form; \leanid{MPSTensor.pairTraceSeparatingAt_threeBlock_of_isNBlkInjective_of_dim_ne}
for the convention-free version.}
as is the conditional directness result: any external proof of the
contradiction $\neg(D_A=D_B\wedge\mathcal G_L^A=\mathcal G_L^B)$ gives
$\mathcal G_{3L}^A\cap\mathcal G_{3L}^B=\{0\}$.\footnote{%
\leanid{MPSTensor.groundSpace_inf_eq_bot_of_not_bondDim_eq_and_groundSpace_eq_of_dim_ge}.}

The equal-size uniqueness branch is isolated as well. For injective tensors,
equal length-$L$ image spaces force equality of the periodic MPV lines. A
supplied distinct-MPV-line hypothesis therefore rules out the equal-size
collapse.\footnote{\leanid{MPSTensor.not_bondDim_eq_and_groundSpace_eq_of_mpvSubmodule_ne}.}
The analytic source of distinct-MPV witnesses is formalized for irreducible
trace-preserving blocks with normalized self-overlaps: if two blocks are not
gauge-phase equivalent, then at arbitrarily large chain length their MPV states
are not pointwise proportional.\footnote{%
\leanid{MPSTensor.exists_ge_not_forall_mpv_eq_mul_of_not_gaugePhaseEquiv_of_irreducible_TP};
\leanid{MPSTensor.exists_ge_not_forall_mpv_eq_mul_of_blocksNotGaugePhaseEquiv_of_irreducible_TP}
for the BNT-family form.}
Combining the uniqueness argument with the finite-dimensional image-space step
gives: if any chain length $N\ge\max\{2,L\}$ witnesses non-proportionality,
then $\mathcal G_{3L}^A\cap\mathcal G_{3L}^B=\{0\}$.\footnote{%
\leanid{MPSTensor.groundSpace_inf_eq_bot_of_exists_not_forall_mpv_eq_mul_of_dim_ge};
\leanid{MPSTensor.groundSpace_inf_eq_bot_of_blocksNotGaugePhaseEquiv_same_dim_of_dim_ge}
for the BNT-separation packaging.}
Zero intersection of image spaces at a fixed length, together with injectivity
of the boundary-to-chain maps at that length, implies homogeneous pair trace
separation.\footnote{\leanid{MPSTensor.pairTraceSeparatingAt_of_groundSpace_inf_eq_bot_of_isNBlkInjective};
and the combined forms
\leanid{MPSTensor.pairTraceSeparatingAt_threeBlock_of_exists_not_forall_mpv_eq_mul_of_dim_ge},
\leanid{MPSTensor.pairTraceSeparatingAt_threeBlock_of_blocksNotGaugePhaseEquiv_same_dim_of_dim_ge}.}
The two branches are combined for all ordered pairs of distinct blocks of a
BNT family at the common length $3L$, and the result is fed into the
selector and product-span construction.\footnote{%
\leanid{MPSTensor.forall_pairTraceSeparatingAt_threeBlock_of_blocksNotGaugePhaseEquiv};
\leanid{MPSTensor.exists_forall_pairTraceSeparatingAt_threeBlock_of_blocksNotGaugePhaseEquiv}.}

The BNT conclusion is stated directly from explicit separated-block
hypotheses, with the required finite-length injectivity hypotheses visible.
These hypotheses may begin at an arbitrary positive common length $L_0$;
positive-length propagation supplies the longer lengths used in the proof.
The selector datum and the resulting finite direct-sum span are theorems rather
than identity-padding assumptions.\footnote{%
\leanid{MPSTensor.hasPairBlockSeparatingWords_threeBlock_of_blocksNotGaugePhaseEquiv};
\leanid{MPSTensor.propBlockInjective_of_blocksNotGaugePhaseEquiv_directSum};
\leanid{MPSTensor.wordTupleSpanTop_of_blocksNotGaugePhaseEquiv_directSum_selectors};
and their Condition-C1 variants in
\path{TNLean/MPS/MPDO/BiCFDerivation/BNTDirectSum.lean}.}

The sharp direct-sum interpolation does not concatenate a separate injective
prefix. Full pair span at a common length $S$ realizes an arbitrary value
$(X,0)$ on a chosen block and one anchor block. Multiplication by $b-2$
identity--zero selectors removes the remaining blocks. Hence the full
simultaneous span holds at length $(b-1)S$. For
$S=3(L_0+1)$ this is exactly the source length
$3(b-1)(L_0+1)$.\footnote{%
\leanid{MPSTensor.wordTupleSpanTop_mul_pred_of_forall_pairWordTupleSpanTop};
\leanid{MPSTensor.wordTupleSpanTop_threeBlock_mul_pred_of_blocksNotGaugePhaseEquiv_c1}.}

The downstream consumers of the direct-sum input are threefold. The exact-length
span of a direct-sum block family and its two-block trace-dual form are expressed
as predicates used throughout the pair route.\footnote{%
\leanid{WordTupleSpanTop}, \leanid{PairWordTupleSpanTop},
\leanid{PairTraceSeparatingAt}.}
Theorems in the period-window module turn a positive homogeneous window of
the pair-span predicate into homogeneous trace separation.\footnote{%
\path{TNLean/MPS/MPDO/BiCFDerivation/PairHomogenization.lean}.}
BNT-level reformulations turn pairwise period-window data into a common
exact-length span for a canonical-form BNT family.\footnote{%
\path{TNLean/MPS/CanonicalForm/BlockDiagonalCommutant.lean}.}

\section{Conclusion}

The arbitrary-$L_0$ direct-sum interpolation needed for the BNT specialization
is now formalized at the exact source length $3(b-1)(L_0+1)$. The resulting
span is homogeneous and positive-length, not cumulative. The length-zero
identity is only the neutral base of a finite product over no remaining
selectors. The case $b=1$ is treated separately by ordinary block injectivity.

For a BNT canonical form of total bond dimension $D$, the common quantum
Wielandt length is $L_0=D^4$. Every representative has positive dimension and
occurs at least once, so $b\leq D$. The formal theorem consequently gives a
positive simultaneous-span length at most
\[
  3(b-1)(D^4+1)\leq3(D-1)(D^4+1)\leq3D^5.
\]
This closes the sharp block-separation input cited by the 2016 MPDO source.\footnote{%
\leanid{MPSTensor.SectorDecomposition.basisCount_le_totalDim};
\leanid{MPSTensor.IsBNTCanonicalForm.exists_basis_wordTupleSpanTop_le_three_totalDim_pow_five}.}

This conclusion concerns the finite-length direct-sum span. Other statements
in the Fundamental Theorem that compare two canonical-form tensors still
require their own reblocking and sector-data arguments; those are logically
separate from the direct-sum input proved here.

\bibliographystyle{alpha}
\bibliography{references}

\end{document}
